% Part: first-order-logic
% Chapter: completeness
% Section: henkin-expansion

\documentclass[../../../include/open-logic-section]{subfiles}

\begin{document}

\olfileid{fol}{com}{hen}

\olsection{हेंकिन विस्तार}

\begin{explain}
संपूर्णता प्रमेय सिद्ध करण्यातील एक आव्हान असे आहे की संपूर्ण सुसंगत
संच~$\Gamma$ पासून आपण रचलेल्या प्रतिमानाने~$\Gamma$ मधील सर्व संख्यकीकृत
!!{formula}s सत्य केली पाहिजेत. याची हमी देण्यासाठी लिऑन हेंकिन यांची एक
युक्ती वापरतो. थोडक्यात, भाषेत अनंत !!{constant}s ची भर घालून तिचा विस्तार
करणे आणि एक मुक्त !!{variable} असलेल्या प्रत्येक !!{formula}~$!A(x)$ साठी
पुढील रूपाचे सूत्र जोडणे, ही ती युक्ती:
\iftag{prvEx}
      {$\lexists[x][!A(x)] \lif !A(c)$}
      {$\lnot\lforall[x][!A(x)] \lif \lnot !A(c)$},
येथे $c$ हे नवीन !!{constant}s पैकी एक आहे. $\Gamma$ पूर्ण करणारी
!!{structure} आपण रचू तेव्हा, यामुळे प्रत्येक
\iftag{prvEx}
{सत्य अस्तित्ववाची वाक्याचा साक्षी}
{असत्य सर्ववाची वाक्याचे प्रतिउदाहरण}
नवीन स्थिरांमध्ये असेल, याची हमी मिळेल.
\end{explain}

\begin{prop}
\ollabel{prop:lang-exp}
$\Gamma$ हा $\Lang L$ मध्ये सुसंगत असेल आणि $\Lang L$ मध्ये नवीन
!!{constant}s चा !!a{denumerable} संच $\Obj d_0$, $\Obj d_1$, \dots
मिळवून $\Lang L'$ तयार केली असेल, तर $\Gamma$ हा~$\Lang L'$ मध्ये
सुसंगत असतो.
\end{prop}

\begin{defn}[संतृप्त संच]
भाषा $\Lang {L}$ मधील !!{formula}s चा संच $\Gamma$
\emph{संतृप्त} असतो तेव्हा आणि केवळ तेव्हाच, जेव्हा एक मुक्त
!!{variable}~$x$ असलेल्या प्रत्येक !!{formula}~$!A(x) \in \Frm[L]$
साठी असे !!a{constant}~$c \in \Lang{L}$ असते की
\iftag{prvEx}
      {$\lexists[x][!A(x)] \lif !A(c) \in \Gamma$}
      {$\lnot\lforall[x][!A(x)] \lif \lnot !A(c) \in \Gamma$}.
\end{defn}

पुढील प्रमेयाच्या सिद्धतेत खालील व्याख्या वापरली जाईल.

\begin{defn}
\ollabel{defn:henkin-exp}
$\Lang L'$ ही \olref{prop:lang-exp} मधीलप्रमाणे असू द्या. ज्या
$\Lang L'$ मधील !!{formula}s~$!A_i(x_i)$ मध्ये एक चर ($x_i$) मुक्तपणे
येतो, त्या सर्वांचे $!A_0(x_0)$, $!A_1(x_1)$, \dots असे प्रगणन निश्चित
करा. आपण~$n$ वर विगमनाने !!{sentence}s~$!D_n$ परिभाषित करतो.

$\Lang{L}$ मध्ये आपण भर घातलेल्या $\Obj d_i$ पैकी~$!A_0(x_0)$ मध्ये
न येणारे पहिले !!{constant} म्हणजे $c_0$ असू द्या. $!D_0$,
\dots,~$!D_{n-1}$ आधीच परिभाषित केली आहेत असे गृहीत धरून, नवीन
!!{constant}s~$\Obj d_i$ पैकी~$!D_0$, \dots,~$!D_{n-1}$ यांमध्ये किंवा
$!A_n(x_n)$ मध्ये न येणारे पहिले स्थिर म्हणजे $c_n$ असू द्या.

आता $!D_{n}$ हे पुढील !!{formula} असू द्या:
\iftag{prvEx}
{$\lexists[x_{n}][!A_{n}(x_{n})] \lif
  !A_{n}(c_{n})$}{$\lnot\lforall[x_{n}][!A_{n}(x_{n})] \lif \lnot
  !A_{n}(c_{n})$}.
\end{defn}

\begin{lem}
\ollabel{lem:henkin}
प्रत्येक सुसंगत संच~$\Gamma$ चा विस्तार करून संतृप्त सुसंगत
संच~$\Gamma'$ मिळवता येतो.
\end{lem}

\begin{proof}
भाषा~$\Lang{L}$ मधील वाक्यांचा सुसंगत संच~$\Gamma$ दिला असता,
नवीन !!{constant}s चा !!a{denumerable} संच जोडून भाषेचा विस्तार
$\Lang{L'}$ पर्यंत करा. \olref{prop:lang-exp} नुसार $\Gamma$ या अधिक
समृद्ध भाषेतही सुसंगत असतो. तसेच $!D_i$ ही
\olref{defn:henkin-exp} मधीलप्रमाणे असू द्या. पुढील व्याख्या करा:
\begin{align*}
\Gamma_0 & = \Gamma \\
\Gamma_{n+1} & = \Gamma_n \cup \{!D_n \}
\end{align*}
म्हणजे $\Gamma_{n+1} = \Gamma \cup \{ !D_0, \dots, !D_n \}$; आणि
$\Gamma' = \bigcup_{n} \Gamma_n$ असू द्या. $\Gamma'$ स्पष्टपणे संतृप्त आहे.

$\Gamma'$ विसंगत असेल, तर कोणत्यातरी $n$ साठी $\Gamma_n$ विसंगत असेल
(सराव: का ते स्पष्ट करा). म्हणून $\Gamma'$ सुसंगत आहे हे दाखवण्यासाठी
प्रत्येक संच~$\Gamma_n$ सुसंगत आहे हे~$n$ वर विगमनाने दाखवणे पुरेसे आहे.

विगमनाचा आधार म्हणजे $\Gamma_0 = \Gamma$ सुसंगत आहे हे विधानच; ते
प्रमेयाचे गृहीतक आहे. विगमन-पायरीसाठी $\Gamma_{n}$ सुसंगत आहे, पण
$\Gamma_{n+1} = \Gamma_n \cup \{!D_n\}$ विसंगत आहे असे समजा. लक्षात
घ्या की $!D_n$ हे
\iftag{prvEx}
{$\lexists[x_{n}][!A_{n}(x_n)] \lif !A_{n}(c_{n})$}
{$\lnot\lforall[x_{n}][!A_{n}(x_n)] \lif \lnot !A_{n}(c_{n})$}
आहे; येथे $!A_n(x_n)$ हे फक्त~$x_n$ चर मुक्त असलेले $\Lang{L'}$ चे
!!a{formula} आहे. $c_n$ निवडण्याच्या पद्धतीनुसार
(\olref{defn:henkin-exp} पाहा), $c_n$ हे~$!A_n(x_n)$ किंवा~$\Gamma_n$
मध्ये येत नाही.

$\Gamma_n \cup \{!D_n\}$ विसंगत असेल, तर $\Gamma_n
\Proves \lnot !D_n$; म्हणून पुढील दोन्ही विधाने खरी आहेत:
\iftag{prvEx}{
\[
\Gamma_n \Proves \lexists[x_n][!A_n(x_n)]
\qquad
\Gamma_n \Proves \lnot !A_n(c_n)
\]}{
\[
\Gamma_n \Proves \lnot\lforall[x_n][!A_n(x_n)]
\qquad
\Gamma_n \Proves !A_n(c_n)
\]}
$c_n$ हे $\Gamma_n$ किंवा~$!A_n(x_n)$ मध्ये येत नसल्यामुळे
\tagrefs{prfAX/{fol:axd:qpr:thm:strong-generalization},prfSC/{fol:seq:qpr:thm:strong-generalization},prfND/{fol:ntd:qpr:thm:strong-generalization},prfTab/{fol:tab:qpr:thm:strong-generalization}}
लागू होते.
\iftag{prvEx}
{$\Gamma_n \Proves \lnot !A_n(c_n)$}
{$\Gamma_n \Proves !A_n(c_n)$}
यापासून अनुक्रमे
\iftag{prvEx}
{$\Gamma_n \Proves \lforall[x_n][\lnot !A_n(x_n)]$}
{$\Gamma_n \Proves \lforall[x_n][!A_n(x_n)]$}
मिळते. अशा रीतीने आपल्याकडे पुढील दोन्ही विधाने आहेत:
\iftag{prvEx}
{$\Gamma_n \Proves \lexists[x_n][!A_n(x_n)]$ आणि
$\Gamma_n \Proves \lforall[x_n][\lnot !A_n(x_n)]$}
{$\Gamma_n \Proves \lnot\lforall[x_n][!A_n(x_n)]$ आणि
$\Gamma_n \Proves \lforall[x_n][!A_n(x_n)]$};
म्हणून $\Gamma_n$ स्वतः विसंगत आहे.
\iftag{prvEx}
{(लक्षात घ्या की
\iftag{prvAll}
{$\lforall[x_n][\lnot !A_n(x_n)] \Proves
  \lnot\lexists[x_n][!A_n(x_n)]$}
{$\lforall[x_n][\lnot !A_n(x_n)]$ ची व्याख्या
  $\lnot\lexists[x_n][\lnot\lnot !A_n(x_n)]$ अशी आहे आणि
  $\lnot\lexists[x_n][\lnot\lnot !A_n(x_n)] \Proves
  \lnot\lexists[x_n][!A_n(x_n)]$}.)}{}
हा विरोधाभास आहे: $\Gamma_n$ सुसंगत असल्याचे गृहीत होते. म्हणून
$\Gamma_n \cup \{ !D_n\}$ सुसंगत आहे.
%READERNOTE{OLCOM-001: गोठवलेल्या इंग्रजी स्रोताच्या व्याख्याविषयक टिपेत सार्वत्रिक सूत्रातील A_n नंतर (x_n) सुटले आहे. मराठीत ते पुनर्स्थापित केले असून गोठवलेले इंग्रजी बाइट्स बदललेले नाहीत.}
\end{proof}

\begin{explain}
आता आपण दाखवू की संतृप्त असलेल्या \emph{संपूर्ण}, सुसंगत संचांना पुढील
गुणधर्म असतो: \iftag{prvAll}{त्यात सर्ववाची संख्यकीकृत !!{sentence}
  तेव्हा आणि केवळ तेव्हाच असते, जेव्हा त्याचे सर्व दृष्टांत त्यात असतात
  }{}\iftag{defAll,defEx}{}{ आणि }\iftag{prvAll}{त्यात अस्तित्ववाची
  संख्यकीकृत !!{sentence} तेव्हा आणि केवळ तेव्हाच असते, जेव्हा त्याचा
  किमान एक दृष्टांत त्यात असतो}{}. संपूर्ण, सुसंगत, संतृप्त संचापासून
आपण निर्माण करू ती !!{structure} त्यातील सर्व संख्यकीकृत वाक्ये सत्य करते,
हे दाखवण्यासाठी हा गुणधर्म वापरू.
\end{explain}

\begin{prop}\ollabel{prop:saturated-instances}
$\Gamma$ हा संपूर्ण, सुसंगत आणि संतृप्त आहे असे समजा.
\begin{tagenumerate}{prvEx,prvAll}
\tagitem{prvEx}{$\lexists[x][!A(x)] \in \Gamma$ तेव्हा आणि केवळ तेव्हाच,
  जेव्हा किमान एका बंद पद~$t$ साठी $!A(t) \in \Gamma$.}{}
\tagitem{prvAll}{$\lforall[x][!A(x)] \in \Gamma$ तेव्हा आणि केवळ तेव्हाच,
  जेव्हा सर्व बंद पदे~$t$ यांच्यासाठी $!A(t) \in \Gamma$.}{}
\end{tagenumerate}
\end{prop}

\begin{proof}
  \begin{tagenumerate}{prvEx,prvAll}
  \tagitem{prvEx}{%
    \iftag{probEx}{सराव.}{प्रथम $\lexists[x][!A(x)]
      \in \Gamma$ असे समजा. $\Gamma$ संतृप्त असल्यामुळे कोणत्यातरी
      !!{constant}~$c$ साठी
      $(\lexists[x][!A(x)] \lif !A(c)) \in \Gamma$. पुढील संदर्भांतील
      \tagrefs{prfAX/{fol:axd:ppr:prop:provability-lif},prfSC/{fol:seq:ppr:prop:provability-lif},prfND/{fol:ntd:ppr:prop:provability-lif},prfTab/{fol:tab:ppr:prop:provability-lif}},
      बाब~(1), आणि
      \olref[ccs]{prop:ccs}\olref[ccs]{prop:ccs-prov-in} नुसार $!A(c)
      \in \Gamma$.

    उलट दिशेसाठी संतृप्तता आवश्यक नाही: $!A(t) \in \Gamma$ असे समजा.
    मग पुढील संदर्भांतील
      \tagrefs{prfAX/{fol:axd:qpr:prop:provability-quantifiers},prfSC/{fol:seq:qpr:prop:provability-quantifiers},prfND/{fol:ntd:qpr:prop:provability-quantifiers},prfTab/{fol:tab:qpr:prop:provability-quantifiers}},
      बाब~(1) नुसार $\Gamma \Proves \lexists[x][!A(x)]$.
    \olref[ccs]{prop:ccs}\olref[ccs]{prop:ccs-prov-in} नुसार
    $\lexists[x][!A(x)] \in \Gamma$.}}{}
  \tagitem{prvAll}{%
    \iftag{probAll}{सराव.}{सर्व बंद पदे~$t$ यांच्यासाठी
      $!A(t) \in \Gamma$ असे समजा. विरोधासाठी
      $\lforall[x][!A(x)] \notin \Gamma$ असे गृहीत धरा. $\Gamma$ संपूर्ण
      असल्यामुळे $\lnot\lforall[x][!A(x)] \in \Gamma$. संतृप्ततेनुसार
      कोणत्यातरी !!{constant}~$c$ साठी
      \iftag{prvEx}{$(\lexists[x][\lnot !A(x)] \lif \lnot !A(c)) \in
        \Gamma$}{$(\lnot\lforall[x][!A(x)] \lif \lnot !A(c)) \in
        \Gamma$}. $c$ हे बंद पद असल्यामुळे गृहीतकानुसार $!A(c) \in
      \Gamma$. पण यामुळे $\Gamma$ विसंगत होईल. (सराव: पुढील संच विसंगत
      असल्याचे दाखवणारी !!{derivation} द्या:
      \[
      \lnot \lforall[x][!A(x)],
      \iftag{prvEx}{\lexists[x][\lnot !A(x)] \lif \lnot
        !A(c)}{\lnot\lforall[x][!A(x)] \lif \lnot
        !A(c)}, !A(c)
      \]
      .)

      उलट दिशेसाठी संतृप्तता आवश्यक नाही:
      $\lforall[x][!A(x)] \in \Gamma$ असे समजा. मग पुढील संदर्भांतील
      \tagrefs{prfAX/{fol:axd:qpr:prop:provability-quantifiers},prfSC/{fol:seq:qpr:prop:provability-quantifiers},prfND/{fol:ntd:qpr:prop:provability-quantifiers},prfTab/{fol:tab:qpr:prop:provability-quantifiers}},
      बाब~(2) नुसार $\Gamma \Proves !A(t)$. \olref[ccs]{prop:ccs}
      नुसार $!A(t) \in \Gamma$.}}{}
  \end{tagenumerate}
\end{proof}

\end{document}
